\vspace{-5pt}
\section{The Conservative Q-Learning (CQL) Framework}
\vspace{-5pt}
In this section, we develop a conservative Q-learning (CQL) algorithm, such that the expected value of a policy under the learned Q-function lower-bounds its true value. A lower bound on the Q-value prevents the over-estimation that is common in offline RL settings due to OOD actions and function approximation error~\citep{levine2020offline,kumar2019stabilizing}. {We use the term CQL to refer broadly to both Q-learning methods and actor-critic methods, though the latter also use an explicit policy.}
We start by focusing on the policy evaluation step in CQL, which can be used by itself as an off-policy evaluation procedure, or integrated into a complete offline RL algorithm, as we will discuss in Section~\ref{sec:framework}.

\vspace{-5pt}
\subsection{Conservative Off-Policy Evaluation}
\label{sec:policy_eval}
\vspace{-5pt}
We aim to estimate the value $V^\policy(\bs)$ of a target policy $\policy$ given access to a dataset, $\mathcal{D}$, generated by following a behavior policy $\behavior(\ba|\bs)$.
Because we are interested in preventing overestimation of the policy value, we learn a \textit{conservative}, lower-bound Q-function by additionally minimizing Q-values alongside a standard Bellman error objective.
Our choice of penalty is to minimize the expected Q-value under a particular distribution of state-action pairs, $\mu(\bs, \ba)$. 
Since standard Q-function training does not query the Q-function value at unobserved states, but queries the Q-function at unseen actions, 
we restrict $\mu$ to match the state-marginal in the dataset,
such that $\mu(\bs, \ba) = d^\behavior(\bs) \mu(\ba|\bs)$.
This gives rise to the iterative update for training the Q-function, as a function of a tradeoff factor $\alpha$:
\begin{equation}
    \small{\hat{Q}^{k+1} \leftarrow \arg\min_{Q}~ \alpha~ \E_{\bs \sim \mathcal{D}, \ba \sim \mu(\ba|\bs)}\left[Q(\bs, \ba)\right] + \frac{1}{2}~ \E_{\bs, \ba \sim \mathcal{D}}\left[\left(Q(\bs, \ba) - \hat{\bellman}^\policy \hat{Q}^{k} (\bs, \ba) \right)^2 \right],} 
    \label{eqn:objective_1}
\end{equation}
In Theorem~\ref{thm:min_q_underestimates}, we show that the resulting Q-function, $\hat{Q}^\policy \coloneqq \lim_{k \rightarrow \infty} \hat{Q}^k$, lower-bounds $Q^\policy$ at all $(\bs,\ba)$. 
However, we can substantially tighten this bound
if we are \textit{only} interested in estimating  $V^\policy(\bs)$. If we only require that the expected value of the $\hat{Q}^\pi$ under $\policy(\ba|\bs)$ lower-bound $V^\pi$, we can improve the bound by introducing an additional Q-value \textit{maximization} term under the data distribution, $\behavior(\ba|\bs)$, resulting in the iterative update (changes from Equation~\ref{eqn:objective_1} in red):
\begin{multline}
    \small{\hat{Q}^{k+1} \leftarrow \arg\min_{Q}~~ \alpha \cdot \left(\E_{\bs \sim \mathcal{D}, \ba \sim \mu(\ba|\bs)}\left[Q(\bs, \ba)\right] - \textcolor{red}{\E_{\bs \sim \mathcal{D}, \ba \sim \hatbehavior(\ba|\bs)}\left[Q(\bs, \ba)\right]} \right)} \\
    \small{+ \frac{1}{2}~ \E_{\bs, \ba, \bs' \sim \mathcal{D}}\left[\left(Q(\bs, \ba) - \hat{\bellman}^\policy \hat{Q}^{k} (\bs, \ba) \right)^2 \right]}.
    \label{eqn:modified_policy_eval}
\end{multline}
In Theorem~\ref{thm:cql_underestimates}, we show that, while the resulting Q-value 
$\hat{Q}^{\policy}$ may not be a point-wise lower-bound, we have $\E_{\policy(\ba|\bs)}[\hat{Q}^\pi(\bs, \ba)] \leq V^\pi(\bs)$
 when $\mu(\ba|\bs) = \policy(\ba|\bs)$. 
Intuitively, since Equation~\ref{eqn:modified_policy_eval} maximizes Q-values under the behavior policy $\hatbehavior$, Q-values for actions that are likely under $\hatbehavior$ might be overestimated, and hence $\hat{Q}^\policy$ may not lower-bound $Q^\policy$ pointwise. 
While in principle the maximization term can utilize other distributions besides $\hatbehavior(\ba|\bs)$, we prove in Appendix~\ref{app:maximizing_distributions} that the resulting value is not guaranteed to be a lower bound for other distribution besides $\hatbehavior(\ba|\bs)$.


\textbf{Theoretical analysis.}
We first note that Equations~\ref{eqn:objective_1} and \ref{eqn:modified_policy_eval} use the empirical Bellman operator, $\hat{\bellman}^\policy$, instead of the actual Bellman operator, $\bellman^\policy$.  Following~\citep{osband2016deep,jaksch2010near,o2018variational}, we use concentration properties of $\hat{\bellman}^\policy$ to control this error. Formally,
for all $\bs, \ba \in \mathcal{D}$, with  probability $\geq 1 - \delta$, $|\hat{\bellman}^\policy - \bellman^\policy|(\bs, \ba) \leq \frac{C_{r,T, \delta}}{\sqrt{|\mathcal{D}(\bs, \ba)|}}$, where $C_{r, T, \delta}$ is a constant dependent on the concentration properties (variance) of $r(\bs, \ba)$ and $\transitions(\bs'|\bs, \ba)$, and $\delta\in (0, 1)$ (see Appendix~\ref{app:handling_unobserved_actions} for more details). 
{For simplicity in the derivation, we assume that $\hatbehavior(\ba|\bs) > 0, \forall \ba \in \mathcal{A}~, \forall \bs \in \mathcal{D}$. $\frac{1}{\sqrt{|\mathcal{D}|}}$ denotes a vector of size $|\mathcal{S}| |\mathcal{A}|$ containing square root inverse counts for each state-action pair, except when $\mathcal{D}(\bs, \ba) = 0$, in which case the corresponding entry is a very large but finite value $\delta \geq \frac{2 R_{\max}}{1 - \gamma}$.}
Now, we show that the conservative Q-function learned by iterating Equation~\ref{eqn:objective_1} lower-bounds the true Q-function. Proofs can be found in Appendix~\ref{app:missing_proofs}.

\begin{theorem} %
\label{thm:min_q_underestimates}
For any $\mu(\ba|\bs)$ with $\supp \mu \subset \supp \hatbehavior$, with probability $\geq 1 - \delta$, $\hat{Q}^\pi$ (the Q-function obtained by iterating Equation~\ref{eqn:objective_1}) satisifies:
\begin{equation*}
\forall \bs \in \mathcal{D}, \ba, ~~ \hat{Q}^\policy(s, a) \leq Q^\policy(\bs, \ba) - \alpha \left[\left(I - \gamma P^\pi \right)^{-1} \frac{\mu}{\hatbehavior} \right](\bs, \ba) + \left[ (I - \gamma P^\pi)^{-1}\frac{C_{r, T, \delta} R_{\max}}{(1- \gamma) \sqrt{|\mathcal{D}|}} \right](\bs, \ba).    
\end{equation*}
Thus, if {$\alpha$ is sufficiently large},
then $\hat{Q}^\policy(\bs, \ba)  \leq Q^\policy(\bs, \ba), \forall \bs \in \mathcal{D}, \ba$. When $\hat{\bellman}^\policy = \bellman^\policy$, any $\alpha > 0$ guarantees $\hat{Q}^\policy(\bs, \ba)  \leq Q^\policy(\bs, \ba), \forall \bs \in \mathcal{D}, \ba \in \mathcal{A}$.
\end{theorem}
Next, we show that Equation~\ref{eqn:modified_policy_eval} lower-bounds the expected value under the policy $\policy$, when $\mu = \policy$. We also show that Equation~\ref{eqn:modified_policy_eval} does not lower-bound the Q-value estimates pointwise. {For this result, we abuse notation and assume that $\frac{1}{\sqrt{|\mathcal{D}|}}$ refers to a vector of inverse square root of only state counts, with a similar correction as before used to handle the entries of this vector at states with zero counts.}
\begin{theorem}[Equation~\ref{eqn:modified_policy_eval} results in a tighter lower bound]
\label{thm:cql_underestimates}
The value of the policy under the Q-function from Equation~\ref{eqn:modified_policy_eval}, $\hat{V}^\policy(\bs) = \E_{\policy(\ba|\bs)}[\hat{Q}^\policy(\bs, \ba)]$, lower-bounds the true value of the policy obtained via exact policy evaluation, $V^\policy(\bs) = \E_{\policy(\ba|\bs)}[Q^\policy(\bs, \ba)]$, when $\mu = \policy$, according to:
\begin{equation*}
\forall \bs \in \mathcal{D}, ~~ \hat{V}^\policy(\bs) \leq V^\policy(\bs) - \alpha \left[\left(I - \gamma P^\pi \right)^{-1} \E_{\policy}\left[\frac{\policy}{\hatbehavior} - 1 \right] \right](\bs) + \left[ (I - \gamma P^\pi)^{-1} \frac{C_{r, T, \delta} R_{\max}}{(1- \gamma) \sqrt{|\mathcal{D}|}}\right](\bs).
\end{equation*}
Thus, if $\alpha > \frac{C_{r, T} R_{\max}}{1 - \gamma} \cdot \max_{\bs \in \mathcal{D}} \frac{1}{|\sqrt{|\mathcal{D}(\bs)|}} \cdot \left[\sum_{\ba} \policy(\ba|\bs) (\frac{\policy(\ba|\bs)}{\hatbehavior(\ba|\bs))} - 1)\right]^{-1}$
,~ $\forall \bs \in \mathcal{D},~\hat{V}^\policy(\bs) \leq {V}^\policy(\bs)$, with probability $\geq 1 - \delta$. When $\hat{\bellman}^\policy = {\bellman}^\policy$, then any $\alpha > 0$ guarantees $\hat{V}^\policy(\bs) \leq V^\policy(\bs), \forall \bs \in \mathcal{D}$.
\end{theorem}
The analysis presented above assumes that no function approximation is used in the Q-function, meaning that each iterate can be represented exactly.
We can further generalize the result in Theorem~\ref{thm:cql_underestimates} to the case of both linear function approximators and non-linear neural network function approximators, where the latter builds on the neural tangent kernel (NTK) framework~\citep{ntk}. Due to space constraints, we present these results in Theorem~\ref{thm:policy_eval_func_approx} and Theorem~\ref{corr:nonlinear_ntk} in Appendix~\ref{app:cql_linear_non_linear}. 

\textbf{In summary}, we showed that the basic CQL evaluation in Equation~\ref{eqn:objective_1} learns a Q-function that lower-bounds the true Q-function $Q^\pi$, and the evaluation in Equation~\ref{eqn:modified_policy_eval} provides a \emph{tighter} lower bound on the expected Q-value of the policy $\pi$. For suitable $\alpha$, both bounds hold under sampling error and function approximation. {We also note that as more data becomes available and $|\mathcal{D}(\bs, \ba)|$ increases, the theoretical value of $\alpha$ that is needed to guarantee a lower bound decreases, which indicates that in the limit of infinite data, a lower bound can be obtained by using extremely small values of $\alpha$.} Next, we will extend on this result into a complete RL algorithm.



